\documentclass[10pt,a4paper]{amsart}
\usepackage[margin=19mm]{geometry}
\usepackage{amsmath,amssymb,mathtools}
\usepackage[scr=rsfs,cal=euler]{mathalfa}
\usepackage{xfrac}
\usepackage[unicode,hidelinks,bookmarks=false]{hyperref}
\newcommand{\ie}{i.e.\ }
\newcommand{\cf}{cf.\ }
\newcommand{\resp}{respectively\ }
\newcommand{\Subsection}{\subsection}
\providecommand{\nobreakdash}{\nobreak\hbox{-}}
\setlength{\emergencystretch}{2em}
\multlinegap0pt
\usepackage{mathtools}
\usepackage{amssymb}
\usepackage{dsfont}
\usepackage{xcolor}
\newtheorem{lem}{Lemma}
\title{\texorpdfstring{Instantaneous Sobolev Regularization \\ for Dissipative Bosonic Dynamics}{Instantaneous Sobolev Regularization for Dissipative Bosonic Dynamics}}
\author{Pablo Costa Rico, Paul Gondolf, Tim M\"{o}bus}
\date{Source: arXiv:2512.04066v2 (CC BY 4.0)}
\begin{document}
\maketitle

\noindent\emph{Source and licence.} \texorpdfstring{Instantaneous Sobolev Regularization \\ for Dissipative Bosonic Dynamics}{Instantaneous Sobolev Regularization for Dissipative Bosonic Dynamics}. Version arXiv:2512.04066v2 (CC BY 4.0), distributed by arXiv under the Creative Commons Attribution 4.0 International licence, http://creativecommons.org/licenses/by/4.0/, as recorded in the arXiv licence metadata for this exact version and shown on the abstract page https://arxiv.org/abs/2512.04066v2. Full licence notice: \texttt{licenses/arxiv-2512.04066-CC-BY-4.0.txt}.\par\smallskip

\noindent\textit{Editorial scope.} Counted unit: the article's lem appx-lem:gronwall-extension (source0.tex lines 1225--1241). The article's complete original proof of it is delivered with it (source0.tex lines 1242--1280, one \texttt{proof} environment of 39 lines). No mathematics is written, repaired or re-derived: the statement and the proof are the article's own lines, in the article's own notation, and no retained line is substituted or re-flowed.  No further source-local context is needed: every cross-reference of every retained line resolves inside the two delivered spans.  Nothing was omitted from any retained span, so the package carries no omission record.  No retained line contains a citation, so the wrapper carries no bibliography and no bibliography entry.  No retained line contains a figure environment, an image-inclusion command or drawing code, so no image is delivered and the package declares no asset dependency.  Editorial formatting only, in addition: an AMS \texttt{amsart} preamble that restates the article's own declarations and macros used by the retained lines, namely the environments \texttt{lem} on separate counters, together with the standard packages those lines need.  The retained environments print in this document as Lemma 1 in order of appearance, whereas the article prints the same items with the numbering of its own full document.  The complete native source of the article as distributed in the arXiv e-print for this exact version is bundled unchanged as \texttt{sources/190307/source0.tex}.
    \begin{lem}\label{appx-lem:gronwall-extension}
        Let $y : [0, \infty) \to [0, \infty)$ be a locally absolutely continuous function satisfying the differential inequality almost everywhere
        \begin{equation*}
            \frac{dy}{dt}(t) \le -a\, y(t)^p + b
        \end{equation*}
        for constants $a>0$, $b\geq0$ and exponent $p > 1$. Then,
        
        \begin{equation*}
            y(t)\leq z(t):=(b/a)^{1/p}+
            \begin{cases}
                ((y(0)-(b/a)^{1/p})^{1-p}+(p-1)at)^{-1/(p-1)},&y(0)>(b/a)^{1/p},\\
                0,&y(0)\leq(b/a)^{1/p}.
            \end{cases}
        \end{equation*}
        

    \end{lem}

    \begin{proof}
        In a first step, we define
        \begin{equation*}
            t_c \coloneqq \inf\{t \ge 0 : y(t) \le (b/a)^{1/p}\}\,.
        \end{equation*}
        For $t > t_c$, it follows immediately that $y(t) \le (b/a)^{1/p}$, and hence $y(t) \le z(t)$. We verify this by contradiction. Suppose there exists $t_v > t_c$ such that $y(t_v) > (b/a)^{1/p}$. By continuity of $y$, the intermediate value theorem implies that there exist $t_0, t_1$ with $t_0 < t_1$ such that
        \begin{equation*}
            y(t_0) = (b/a)^{1/p}, \qquad y(t_1) > (b/a)^{1/p}, \qquad \text{and} \quad y(t) \ge (b/a)^{1/p} \quad\text{ for all } t \in [t_0, t_1]\,.
        \end{equation*}
        
        Since $y(t)\geq(b/a)^{1/p}$ on $[t_0,t_1]$, the differential inequality gives $y'(t)\leq0$ almost everywhere on this interval. Absolute continuity therefore implies
        \begin{equation*}
            y(t_1)-y(t_0)=\int_{t_0}^{t_1}y'(s)\,ds\leq0,
        \end{equation*}
        
        a contradiction. Therefore, $y(t) \le (b/a)^{1/p}$ for all $t \ge t_c$.

        For $0 \le t < t_c$, define $g(t) \coloneqq y(t) - (b/a)^{1/p} > 0$. Then $g'(t) = y'(t)$ almost everywhere, and the differential inequality yields
        \begin{equation*}
            g'(t) \le -a\, g(t)^p.
        \end{equation*}
        Separating variables and integrating from $0$ to $t$ gives
        \begin{equation*}
            \frac{g(t)^{1 - p} - g(0)^{1 - p}}{1 - p} = \int_0^t \frac{g'(s)}{g(s)^p}\, ds \le \int_0^t -a\, ds = -a t.
        \end{equation*}
        Rearranging terms, we obtain
        \begin{equation*}
            g(t)^{1 - p} \ge g(0)^{1 - p} + (p - 1)a t.
        \end{equation*}
        Taking both sides to the power of $\tfrac{1}{1 - p}$ (which reverses the inequality since $\tfrac{1}{1 - p} < 0$), we find
        \begin{equation*}
            g(t) \le \frac{1}{\big(g(0)^{1 - p} + (p - 1)a t \big)^{\frac{1}{p - 1}}}.
        \end{equation*}
        Substituting back $g(t) = y(t) - (b/a)^{1/p}$ yields
        \begin{equation*}
            y(t) \le \frac{1}{\big((y(0) - (b/a)^{1/p})^{1 - p} + (p - 1)a t \big)^{\frac{1}{p - 1}}} + (b/a)^{1/p}.
        \end{equation*}
        Together with the case $t\geq t_c$, this completes the proof, including $b=0$.
    \end{proof}

\end{document}
